\section{Experiments}
\label{sec:experiments}

\subsection{Experimental setup}

\paragraph{Datasets.}
We use the Industrial and Office Amazon benchmarks already prepared in the MiniOneRec reproduction pipeline. Industrial contains $36{,}259$ training interactions, $4{,}533$ test interactions, and $3{,}686$ items. Office contains $38{,}924$ training interactions, $4{,}866$ test interactions, and $3{,}459$ items. Both use the repository's existing time-based splits.

\paragraph{Tokenizer backbone.}
Unless otherwise noted, all training-time tokenizer comparisons use the upstream-aligned MiniOneRec \rqvae{} regime established in \Cref{tab:provenance}: three quantizers with codebook size $256$, embedding dimension $32$, AdamW, learning rate $10^{-3}$, $10{,}000$ epochs, and batch size $20{,}480$.

\paragraph{Evaluation metrics.}
We report three classes of metrics.
\begin{itemize}[leftmargin=1.4em]
  \item \textbf{Probe metrics}: target-better rate on ambiguity buckets, especially same-$l_2$.
  \item \textbf{Tokenizer metrics}: best train-stage collision, final \texttt{sid-generate} collision, and maximum conflict size.
  \item \textbf{Local ambiguity metrics}: mean target $l_2$ leaf count, the fraction of targets in multi-leaf same-$l_2$ buckets, the fraction of targets in deep crowded same-$l_2$ buckets, and conditional entropy of level 3 under a fixed level-1/2 prefix.
\end{itemize}

\subsection{Design-space probes: spectral middle views matter}

We began with a non-invasive probe protocol that does not modify the tokenizer. Given an existing top-1 SID error, we ask whether a graph view scores the ground-truth target item above the currently favored wrong candidate. This lets us study which graph view is most informative before injecting it into training.

\begin{table}[t]
  \centering
  \caption{Best middle-view probe results on same-$l_2$ buckets. Spectral middle views are much stronger than the first heuristic middle views, which motivated the use of a spectral $G_{\mathrm{mid}}$ in MGR-SID v1.}
  \label{tab:probe}
  \begin{tabular}{lccc}
    \toprule
    Dataset & Best heuristic middle view & Best spectral middle view & Spectral view name \\
    \midrule
    Industrial & 0.2747 & 0.5401 & \texttt{fagsp\_mid\_base} \\
    Office & 0.2381 & 0.7619 & \texttt{fagsp\_mid\_base} \\
    \bottomrule
  \end{tabular}
\end{table}

\Cref{tab:probe} provides the first direct justification for the graph bank design. The diagnostic result is not just that graph signals help. Rather, the result is that \emph{middle-resolution graph structure matters enough to deserve explicit design}. The same-$l_2$ gain from the best spectral middle view is nearly doubled on Industrial and more than tripled on Office relative to the first heuristic middle view. This is what led us to instantiate $G_{\mathrm{mid}}$ with a spectral graph rather than a crude diffusion proxy.

\subsection{Training-time comparison under upstream-aligned tokenizer training}

We next test whether the graph signal survives once it is injected into tokenizer learning. All variants here use the fair upstream-aligned training setup from \Cref{tab:provenance}. We compare the semantic baseline, the uniform regularization control, and the hierarchy-aware MGR-SID v1 regularizer.

\begin{table}[t]
  \centering
  \caption{Industrial tokenizer results under the upstream-aligned MiniOneRec training regime. Train-stage collision alone is misleading: the hierarchy-aware variant becomes the best tokenizer after \texttt{sid-generate}.}
  \label{tab:main}
  \begin{tabular}{lccc}
    \toprule
    Method & Best train collision & Final index collision & Max conflict \\
    \midrule
    Semantic baseline & 0.1066 & 0.00353 & 2 \\
    Uniform graph regularization & 0.2458 & 0.00597 & 6 \\
    Hierarchy-aware MGR-SID v1 & 0.1319 & \textbf{0.00326} & 2 \\
    \bottomrule
  \end{tabular}
\end{table}

\Cref{tab:main} reveals two important facts. First, uniform graph regularization is clearly harmful. This matters because it shows that adding graph structure per se is not enough. Second, train-stage collision and final tokenizer quality are not aligned. The semantic baseline still has the best raw train-stage collision, but after the official \texttt{sid-generate} stage, the hierarchy-aware tokenizer becomes the best final index. This is exactly why our evaluation protocol treats the final exported SID as the primary tokenizer artifact.

\subsection{Does the final tokenizer actually reduce local ambiguity?}

The main question is whether the better final collision reflects the failure mode we care about. To answer this, we compare the final baseline index and the final hierarchy-aware index at the level of same-$l_2$ crowding.

\begin{table}[t]
  \centering
  \caption{Local ambiguity comparison on Industrial final indices. The hierarchy-aware tokenizer reduces both catalog-level and test-weighted same-$l_2$ crowding.}
  \label{tab:ambiguity}
  \begin{tabular}{lccc}
    \toprule
    Metric & Baseline & Hierarchy & Delta \\
    \midrule
    Mean target $l_2$ leaf count & 4.7999 & 4.4498 & -0.3501 \\
    Fraction targets in multi-leaf $l_2$ & 0.6894 & 0.6131 & -0.0763 \\
    Fraction targets in $l_2$ with $\ge 4$ leaves & 0.3283 & 0.2828 & -0.0454 \\
    Mean target level-3 entropy under $l_2$ & 1.4533 & 1.2935 & -0.1598 \\
    \bottomrule
  \end{tabular}
\end{table}

\Cref{tab:ambiguity} shows that the improvement is structurally meaningful. The hierarchy-aware tokenizer does not merely shuffle code assignments while preserving the same ambiguity. It lowers the average number of competing leaves under the target prefix, reduces the fraction of targets that fall into multi-leaf same-$l_2$ buckets, and lowers the conditional entropy of the third-level code. In the full movement analysis, $17.32\%$ of test targets move from a multi-leaf same-$l_2$ bucket into a singleton bucket, while only $9.68\%$ move in the opposite direction. This strongly supports our original diagnosis that the remaining tokenizer bottleneck is local leaf ambiguity.

\subsection{First downstream test: v1 transfers only partially}

We first pushed the final Industrial tokenizers from \Cref{tab:main,tab:ambiguity} into downstream SFT. Here, the internal comparison uses the upstream-aligned semantic run and the hierarchy-aware v1 run under the same controlled pipeline. We emphasize that this \emph{internal control} is useful for diagnosis, but it is not our final baseline anchor.

\begin{table}[t]
  \centering
  \caption{Industrial SFT results for the internal control comparison. MGR-SID v1 improves short-range ranking but does not yet form a broad $@10+$ gain.}
  \label{tab:sft_v1}
  \begin{tabular}{lcccccc}
    \toprule
    Run & N@1 & N@3 & N@10 & H@1 & H@3 & H@10 \\
    \midrule
    Internal semantic control & 0.0631 & 0.0777 & 0.0943 & 0.0631 & 0.0882 & 0.1343 \\
    Internal hierarchy-aware v1 & 0.0627 & 0.0802 & 0.0936 & 0.0627 & 0.0929 & 0.1304 \\
    \bottomrule
  \end{tabular}
\end{table}

\Cref{tab:sft_v1} was the first indication that tokenizer gains do transfer downstream, but only partially. Relative to the internal semantic control, v1 improves $\mathrm{NDCG}@3$ and $\mathrm{HR}@3$, and slightly improves $@5$, but remains weaker at $@10+$. This pattern motivated the next step: instead of further changing the graph carrier, we redesigned the training interface itself.

\subsection{Ambiguity-aware tokenizer v2}

The top-$k$ and error-distribution analysis of v1 suggested that graph regularization was helpful on crowded local ambiguity cases but over-corrected already-stable semantic regions. We therefore introduced an ambiguity-aware v2 tokenizer with two changes:
\begin{itemize}[leftmargin=1.4em]
  \item an \emph{offline ambiguity prior} built from semantic density, semantic-collaborative disagreement, and graph competition strength;
  \item a \emph{semantic-structure retention} regularizer that preserves the original semantic neighborhood in low-ambiguity regions.
\end{itemize}

The first v2 tokenizer run used only the offline combined ambiguity prior, because the current online uncertainty proxy was not reliable enough in preliminary sanity checks.

\begin{table}[t]
  \centering
  \caption{Tokenizer-side comparison after introducing ambiguity-aware supervision. The v2 tokenizer does not further reduce final collision, but produces the cleanest local structure.}
  \label{tab:v2tok}
  \begin{tabular}{lcccc}
    \toprule
    Method & Final collision & Mean target $l_2$ leaves & Multi-leaf same-$l_2$ & Mean level-3 entropy \\
    \midrule
    Internal semantic control & 0.00353 & 4.7999 & 0.6894 & 1.4533 \\
    Hierarchy-aware v1 & \textbf{0.00326} & 4.4498 & 0.6131 & 1.2935 \\
    Ambiguity-aware v2 & 0.00353 & \textbf{4.3422} & \textbf{0.4873} & \textbf{1.1001} \\
    \bottomrule
  \end{tabular}
\end{table}

\Cref{tab:v2tok} is important for two reasons. First, it shows that collision is not the whole tokenizer story: v2 only ties the internal semantic control on final collision, yet it yields the cleanest local hierarchy among all tested tokenizers. Second, the ambiguity-aware design does exactly what it was intended to do: it reduces the fraction of targets that remain trapped in multi-leaf same-$l_2$ buckets without simply increasing graph pressure everywhere.

\subsection{Downstream SFT under the recipe-aligned original MiniOneRec setting}

We first evaluate whether this cleaner v2 tokenizer helps in a fair downstream comparison against the \emph{original} MiniOneRec under the same task recipe. The first v2 downstream run uses:
\begin{itemize}[leftmargin=1.4em]
  \item \texttt{title\_history2sid = on}
  \item \texttt{SID-description alignment = off}
\end{itemize}
Hence the recipe-aligned original MiniOneRec comparison is the historical run with the same task setting. For context, we also report the strongest original MiniOneRec SFT and RL results in the repository.

\begin{table}[t]
  \centering
  \caption{Industrial downstream comparison for v2. V2 already surpasses the recipe-aligned original MiniOneRec on all reported NDCG cutoffs and on HR@1/3/5, but does not yet surpass the strongest original MiniOneRec recipe.}
  \label{tab:v2sft}
  \begin{tabular}{lcccccc}
    \toprule
    Run & N@1 & N@3 & N@10 & H@1 & H@3 & H@10 \\
    \midrule
    Recipe-aligned original MiniOneRec & 0.0624 & 0.0798 & 0.0987 & 0.0624 & 0.0929 & \textbf{0.1471} \\
    Strongest original MiniOneRec SFT & 0.0671 & \textbf{0.0850} & \textbf{0.1037} & 0.0671 & \textbf{0.0984} & 0.1509 \\
    Strongest original MiniOneRec RL & \textbf{0.0732} & 0.0890 & 0.1073 & \textbf{0.0732} & 0.1004 & 0.1513 \\
    Ambiguity-aware v2 & 0.0704 & 0.0839 & 0.1008 & 0.0704 & 0.0942 & 0.1425 \\
    \bottomrule
  \end{tabular}
\end{table}

\Cref{tab:v2sft} establishes the current position of our method much more clearly than the earlier internal control. Relative to the \emph{recipe-aligned original MiniOneRec}, v2 improves all reported NDCG cutoffs:
\[
\Delta \mathrm{NDCG}@1 = +0.0079,\;
\Delta \mathrm{NDCG}@3 = +0.0041,\;
\Delta \mathrm{NDCG}@5 = +0.0036,\;
\Delta \mathrm{NDCG}@10 = +0.0021.
\]
It also improves $\mathrm{HR}@1$, $\mathrm{HR}@3$, and $\mathrm{HR}@5$, while trailing slightly on $\mathrm{HR}@10$.

This means the ambiguity-aware tokenizer is no longer just ``working on an internal control.'' It already improves downstream recommendation quality against the original MiniOneRec under the same task recipe. At this stage, however, it still does \emph{not} surpass the strongest original MiniOneRec recipe, which disables \texttt{title\_history2sid} and enables description alignment with probability $0.5$.

\subsection{Recipe isolation: the main mismatch is \texttt{title\_history2sid\_off}}

To understand why the strongest original recipe did not transfer cleanly to v2, we completed the missing cells of the v2 recipe matrix on Industrial:
\begin{itemize}[leftmargin=1.4em]
  \item \texttt{title\_history2sid\_on + desc\_align\_off}
  \item \texttt{title\_history2sid\_on + desc\_align\_p05}
  \item \texttt{title\_history2sid\_off + desc\_align\_off}
  \item \texttt{title\_history2sid\_off + desc\_align\_p05}
\end{itemize}

\begin{table}[t]
  \centering
  \caption{Recipe isolation for the v2 tokenizer on Industrial. The best downstream recipe for v2 is \texttt{title\_history2sid\_on + desc\_align\_p05}, and the dominant negative factor is turning \texttt{title\_history2sid} off.}
  \label{tab:v2recipe}
  \begin{tabular}{lcccccc}
    \toprule
    Run & N@1 & N@3 & N@10 & H@1 & H@3 & H@10 \\
    \midrule
    \texttt{v2\_on\_off} & 0.0704 & 0.0839 & 0.1008 & 0.0704 & 0.0942 & 0.1425 \\
    \texttt{v2\_on\_p05} & \textbf{0.0706} & \textbf{0.0845} & \textbf{0.1027} & \textbf{0.0706} & \textbf{0.0951} & \textbf{0.1463} \\
    \texttt{v2\_off\_off} & 0.0589 & 0.0741 & 0.0913 & 0.0589 & 0.0854 & 0.1339 \\
    \texttt{v2\_off\_p05} & 0.0591 & 0.0737 & 0.0899 & 0.0591 & 0.0843 & 0.1308 \\
    \bottomrule
  \end{tabular}
\end{table}

\Cref{tab:v2recipe} resolves the strongest-recipe mismatch. When \texttt{title\_history2sid} stays on, adding description alignment with probability $0.5$ is beneficial:
\[
\Delta \mathrm{NDCG}@10 = +0.0019,\qquad
\Delta \mathrm{HR}@10 = +0.0038.
\]
By contrast, turning \texttt{title\_history2sid} off causes a large drop regardless of whether description alignment is enabled. Under \texttt{desc\_align\_off}, the drop is
\[
\Delta \mathrm{NDCG}@10 = -0.0096,\qquad
\Delta \mathrm{HR}@10 = -0.0086,
\]
and under \texttt{desc\_align\_p05} it becomes even larger:
\[
\Delta \mathrm{NDCG}@10 = -0.0128,\qquad
\Delta \mathrm{HR}@10 = -0.0154.
\]

This means the main negative factor is not description alignment itself. The real source of mismatch is \texttt{title\_history2sid\_off}. For the current v2 tokenizer, the best downstream setting is therefore \texttt{title\_history2sid\_on + desc\_align\_p05}.

We also compared this best v2 setting against the strongest original MiniOneRec SFT. The remaining gap is now very small:
\[
\Delta \mathrm{NDCG}@10 = -0.0010,\qquad
\Delta \mathrm{HR}@10 = -0.0046.
\]
Further top-$k$ analysis shows that v2 already wins at rank 1 and produces much cleaner same-prefix error behavior, while the strongest original SFT still has an advantage in mid-beam neighborhood retention at \texttt{top3/top5/top10}. In other words, the remaining issue is no longer a tokenizer-collapse problem; it is a downstream interaction gap.

\subsection{Current interpretation}

Taken together, the experiments support a more precise claim than our earlier v1-only narrative:
\begin{enumerate}[leftmargin=1.4em]
  \item fair tokenizer comparison requires the upstream MiniOneRec training regime;
  \item graph-structured collaborative information has real signal for local same-$l_2$ disambiguation, and this signal is strongest when encoded by a spectral middle-resolution graph;
  \item uniform hierarchy-aware graph injection is insufficient because it over-corrects already-stable semantic regions;
  \item ambiguity-aware graph supervision plus semantic-structure retention yields a tokenizer with a cleaner local hierarchy than both the internal semantic control and v1;
  \item this cleaner tokenizer already improves downstream recommendation against the recipe-aligned original MiniOneRec;
  \item after recipe isolation, the best downstream setting for the current v2 tokenizer is \texttt{title\_history2sid\_on + desc\_align\_p05};
  \item the strongest original MiniOneRec recipe remains slightly better overall, but the remaining gap is now small and is best described as a mid-beam neighborhood-retention gap rather than a tokenizer-side failure.
\end{enumerate}

At the current stage, the main conclusion is therefore not ``we already beat the whole system.'' The stronger conclusion is:
\begin{quote}
the tokenizer-side method is now supported by structural evidence, recipe-aligned downstream evidence, and recipe-isolation evidence. The current best v2 setting is already very close to the strongest original MiniOneRec SFT, and the remaining gap is primarily a downstream retention issue rather than a failure of the ambiguity-aware tokenizer itself.
\end{quote}
